\subsection{I-39. Two exact failed repairs for unrestricted mixed affine words}
\label{I-39}

\paragraph{Scope.}
We prove two obstructions to the next proposed extensions of I-38.
First, the ordinary common-center one-pair map can lower a strict
internal interval from generic rank two to rank one even in an
irreducible multilinear width-three scalar word whose every generic
interval rank is two or three. Second, the symmetric semiseparable
Hessian pattern, all three tangent kernel identities, and the whole
additional-pole quadratic identity do not by themselves force the
Hessian to be zero. Neither construction is a counterexample to
the centered two-pair sampler. The first scalar is hit already
by one pair; the second is an abstract family of quadratic jets,
with a stated failure of integration to a fixed polynomial.
Both constructions are over characteristic zero.

\paragraph{1. A rational identity and its residue forms.}
We recall, and verify the needed parts of, the explicit construction
of P-16. Put $a=\sqrt{3/2}$ and, for an integer $d\ge2$, set
\[
 \begin{aligned}
 P_1(T)&=T^{2d}-1,&P_2(T)&=T^d-a,\\
 P_3(T)&=-2(T^d+a),&P_4(T)&=T^{2d}-2,\\
 R_1&=P_2/P_1,&R_2&=1/P_2,\\
 R_3&=1/P_3,&R_4&=P_3/P_4.
 \end{aligned}
 \tag{I.107}
\]
Work in a splitting field of these polynomials. Their four pole
sets are disjoint and simple, of sizes $2d,d,d,2d$ respectively;
none of the displayed numerators cancels a pole. For a pole
$b$ of $R_j$, let $\rho_b\ne0$ be its residue. Partial fractions
and decay $R_j(T)=O(T^{-d})$ give
\[
 R_j(T)=\sum_{b\in\mathcal A_j}\frac{\rho_b}{T-b},
 \qquad \sum_{b\in\mathcal A_j}\rho_b b^e=0
                    \quad(0\le e\le d-2).
 \tag{I.108}
\]
Moreover $P_2P_3=-P_1-P_4$, so
$R_1R_2+R_1R_4+R_3R_4=0$.
For the disjointly supported nonzero linear forms
$L_j=\sum_{b\in\mathcal A_j}\rho_b X_b$, define
\[
 Q=L_1L_2+L_1L_4+L_3L_4\ne0.
 \tag{I.109}
\]
Nonvanishing follows, for example, from any coefficient using
one variable from group one and one from group two. If $g$ has
degree at most $d-1$, division by $T-u$ and (I.108) give
\[
 L_j\left(\frac{g(b)}{b-u}:b\in\bigcup\mathcal A_j\right)
 =-g(u)R_j(u).
 \tag{I.110}
\]
Therefore $Q$ vanishes on all these numerator-at-most-$(d-1)$,
one-pole evaluations. These are the only P-16 facts used below.

\paragraph{2. Ordinary one-pair rank loss on an irreducible affine word.}
Take $d=2$, giving twelve core nodes. They come in opposite pairs,
and every $R_j$ is an even rational function. Add the five nodes
$0,3,-3,4,-4$, giving seventeen distinct nodes in total. Assign
the four nonzero new nodes to variables $s,u,v,t$, respectively;
the coordinate at zero is unused by the input polynomial.
Let
\[
 A(T)=\prod_{b\text{ among all seventeen nodes}}(T-b),
 \qquad d_b=A'(b),\qquad
 L_j^*=\sum_{b\in\mathcal A_j}\rho_b d_bX_b.
 \tag{I.111}
\]
The polynomial $A$ is odd and $A'$ is even. The residues of an
even rational function at opposite poles are negatives of one
another. It follows that $\sum_{b\in\mathcal A_j}\rho_b d_b=0$
for every group. Hence under the ordinary common-center map
$X_b=\tau+z\lambda_b(y)$ on these seventeen nodes,
\[
 L_j^*=zA(y)R_j(y),\qquad
 Q^*=L_1^*L_2^*+L_1^*L_4^*+L_3^*L_4^*=0.
 \tag{I.112}
\]
Indeed $d_b\lambda_b(y)=A(y)/(y-b)$, so this is an exact
polynomial identity after clearing the displayed rational notation.
All coefficients $\rho_bd_b$ are nonzero, and $Q^*$ is a
nonzero quadratic before substitution.

Put $A_0=L_1^*+L_3^*$, $B_0=L_2^*+L_4^*$ and $C_0=L_2^*L_3^*$.
The middle word consists of the affine unipotent factors
$I-\rho_bd_bX_bE_{21}$ in groups one and three, and
$I-\rho_bd_bX_bE_{32}$ in groups two and four, in group order.
Within each group their product is $I-L_j^*E$, since $E^2=0$.
Thus its full product is exactly
\[
 M=I-A_0E_{21}-B_0E_{32}+C_0E_{31}.
 \tag{I.113}
\]
Let $P_L=\operatorname{diag}(0,1,1)$ and
$P_R=\operatorname{diag}(1,1,0)$. Their corner is
\[
 P_LMP_R=
 \begin{pmatrix}0&0&0\\-A_0&1&0\\C_0&-B_0&0\end{pmatrix},
 \qquad
 \det\begin{pmatrix}-A_0&1\\C_0&-B_0\end{pmatrix}=Q^*.
 \tag{I.114}
\]
Use genuinely variable-dependent singular anchor factors
\[
 J_L(u)=(I+uE_{32})P_L,\qquad
 J_R(v)=P_R(I+vE_{21}).
\]
They have rank two at every value. The strict internal interval
$J_L(u)M J_R(v)$ has generic rank two because $Q^*\ne0$.
It has rank exactly one at the center (I.112): its indicated
two-by-two determinant is zero, while its $(2,2)$ entry remains
one. The outer unipotent multipliers in the anchors preserve
this rank calculation.

Now take the actual scalar word
\[
 F=e_3^{\mathsf T}(I+sE_{32})J_L(u)M J_R(v)
                      (I+tE_{21})e_1
   =(s+u-B_0)(t+v-A_0)-Q^*.
 \tag{I.115}
\]
Every variable is read once; every matrix factor is affine.
The two anchors are the only singular factors. An interval
containing no anchor is invertible, one containing exactly one
anchor has rank two, and one containing both has rank two
because (I.114) has nonzero determinant. Thus all generic
interval ranks are in $\{2,3\}$.

The scalar $F$ is irreducible. The invertible polynomial change
of variables $S=s+u-B_0$, $T=t+v-A_0$, keeping the remaining
variables fixed, converts it to $ST-Q^*$. As a polynomial in
$T$ over the UFD of the other variables, its coefficients $S$
and $-Q^*$ are coprime. It is primitive and linear, hence
irreducible over that UFD by Gauss's lemma. This also proves
irreducibility after extending to the algebraic closure.

Nevertheless $F$ is detected by the very center that loses
the internal rank. By (I.112), its image is the product of
the two factors in (I.115). Each has coefficient two on
$\tau$, while $A_0,B_0$ have no $\tau$ term. The image has
coefficient four on $\tau^2$ and is nonzero. This demonstrates
that irreducibility alone does not repair one-pole interval-rank
preservation, and that requiring preservation of every interval
rank is stronger than the desired scalar nonvanishing.

\paragraph{3. The longer-block Hessian conditions alone also fail.}
Take $d=4$ in part~1, with its twenty-four core nodes, and let
$H_0$ be the symmetric zero-diagonal Hessian of $Q$ in (I.109).
In the group order, define vectors $\alpha,\beta$ by putting
$\alpha_b=\rho_b$ in groups one and three and zero otherwise,
and $\beta_b=\rho_b$ in groups two and four and zero otherwise.
For $i<j$ the upper entries satisfy
$(H_0)_{ij}=\alpha_i\beta_j$. This is exactly the single-block
upper rank-one semiseparable pattern arising from an arbitrary
invertible stretch in the I-38 Hessian calculation.

Let $h_i(y)=1/(a_i-y)$, where $a_i$ now lists the twenty-four
core nodes, and set
\[
 D(y)=\operatorname{diag}((a_i-y)^2),\qquad
 H(y)=D(y)H_0D(y).
 \tag{I.116}
\]
It is nonzero, symmetric and zero-diagonal, and its upper
entries still factor as
$[\alpha_i(a_i-y)^2][\beta_j(a_j-y)^2]$.
Every $L_j$ kills polynomial evaluations of degrees at most
two by (I.108). Therefore $H_0$ annihilates each such evaluation
vector, and
\[
 H(y)\mathbf1=H(y)\mathbf h(y)
              =H(y)\mathbf h(y)^{\circ2}=0.
 \tag{I.117}
\]
The three vectors entering $H_0$ here are respectively
$((a_i-y)^2)_i$, $((a_i-y))_i$, and $\mathbf1$.

It also satisfies the complete extra-pole quadratic identity:
for an independent pole $u$,
\[
 \mathbf h(u)^{\mathsf T}H(y)\mathbf h(u)
 =2Q\left(\frac{(a_i-y)^2}{a_i-u}:1\le i\le24\right)=0.
 \tag{I.118}
\]
This is (I.110) with numerator $(T-y)^2$ of degree two.
Consequently the tangent kernels and the additional-pole
quadratic restriction do not suffice, even with the exact
semiseparable pattern, to force $H=0$.

This family is not itself the Hessian pullback of a fixed
polynomial at $\tau\mathbf1+z\mathbf h(y)$. If it were,
specializing $z=0$ would make that pullback independent of $y$,
whereas every nonzero entry of (I.116) is a nonconstant
degree-four polynomial in $y$. Thus the construction explicitly
exhibits the missing coordinate-ring and Hessian-integrability
constraints; it does not silently assert they have been met.
Multiplying or clearing denominators to enter the center's
coordinate ring would still require checking those constraints,
and the original width-three affine representation, anew.
\hfill$\square$

\paragraph{Provenance and next gate.}
The rational identity and residue forms are from the fully
checked P-16 construction; all extra matrix, parity, irreducibility
and Hessian deductions needed here are written out. No historical
priority claim is made. The remaining two-pair target must use
the assumed scalar annihilation, rather than irreducibility
alone, to control rank loss; or retain the integrability and
affine compound constraints beyond the displayed Hessian
identities. Neither obstruction settles that target or the
growing-width polynomial-size problem.
